\documentclass[11pt]{article}

\usepackage[margin=0.85in]{geometry}
\usepackage{amsmath,amssymb}
\usepackage{enumitem}
\usepackage{xcolor}
\usepackage{tcolorbox}
\usepackage{fancyhdr}
\usepackage{lastpage}
\usepackage{hyperref}
\usepackage{microtype}

\definecolor{uwpurple}{HTML}{4B2E83}
\definecolor{uwgold}{HTML}{B7A57A}
\definecolor{softpurple}{HTML}{F4F1FA}

\hypersetup{
    colorlinks=true,
    linkcolor=uwpurple,
    urlcolor=uwpurple
}

\pagestyle{fancy}
\fancyhf{}
\lhead{\textbf{MATH 394: Probability I}}
\rhead{\textbf{Homework}}
\cfoot{\thepage\ of \pageref{LastPage}}

\setlength{\parindent}{0pt}
\setlength{\parskip}{0.45em}

\newcommand{\course}{MATH 394: Probability I}
\newcommand{\instructor}{Arman Jahangiri}
\newcommand{\department}{Department of Mathematics}
\newcommand{\university}{University of Washington}

\newtcolorbox{problemheader}{
    colback=softpurple,
    colframe=uwpurple,
    boxrule=0.8pt,
    arc=2mm,
    left=2mm,
    right=2mm,
    top=1mm,
    bottom=1mm
}

\newcommand{\source}[1]{%
    \vspace{0.15cm}
    {\small\textit{Source: Blitzstein and Hwang, \textit{Introduction to Probability}, Chapter 1, Problem #1.}}
}
\newcommand{\answerbox}{%
\begin{flushright}
\begin{tcolorbox}[
    width=0.36\textwidth,
    height=2cm,
    colback=white,
    colframe=uwpurple,
    boxrule=0.7pt,
    arc=1mm
]
{\small\textbf{Final Answer}}
\end{tcolorbox}
\end{flushright}
}

\newcommand{\partanswerbox}[1]{%
\begin{flushright}
\begin{tcolorbox}[
    width=0.36\textwidth,
    height=2cm,
    colback=white,
    colframe=uwpurple,
    boxrule=0.7pt,
    arc=1mm
]
{\small\textbf{Final Answer for Part #1}}
\end{tcolorbox}
\end{flushright}
}

\begin{document}

\begin{titlepage}
\centering
\vspace*{1.2cm}

{\Large \textbf{\university}}\\
{\large \department}

\vspace{1.5cm}

{\Huge \color{uwpurple}\textbf{Homework 1}}\\[0.25cm]
{\Large \textbf{\course}}\\[0.25cm]
{\large Instructor: Arman Jahangiri}

\vspace{1cm}

\begin{tcolorbox}[
    width=0.78\textwidth,
    colback=softpurple,
    colframe=uwpurple,
    boxrule=0.9pt,
    arc=3mm
]
\centering

\textbf{Submission:} Single PDF on Gradescope

\textbf{Deadline:} 10:00 PM Pacific Time on the listed due date
\end{tcolorbox}

\vspace{0.8cm}

\begin{tcolorbox}[
    width=0.86\textwidth,
    colback=white,
    colframe=uwgold,
    boxrule=0.8pt,
    arc=3mm
]
This homework is worth \textbf{50 points}.

Across the quarter, there are \textbf{eight homework assignments}, worth \textbf{400 points total}. Homework assignments together account for \textbf{40\%} of the final course grade. Further course policies are can be seen in the \href{https://armanjg.github.io/MATH-394-Probability-1-v2/}{\textbf{MATH 394 syllabus}.}
\end{tcolorbox}

\vfill


\end{titlepage}

\section*{Homework Policy}

\subsection*{Submission and Deadline}

Please:

\begin{itemize}[leftmargin=*]
    \item Submit homework as a single PDF on Gradescope by \textbf{10:00 PM (Pacific Time)} on the listed due date.
    \item Begin each problem on a new page, and ensure that pages are correctly tagged in Gradescope.
    \item Write solutions clearly and legibly.
    \item Typesetting solutions in LaTeX is encouraged, but not required.
\end{itemize}

\subsection*{Late Days}

Each student is allotted \textbf{six late days} for the quarter. A late day extends a homework deadline by up to 24 hours without penalty. For example, submitting an assignment anytime between 10:01 PM on the due date and 10:00 PM the following day counts as one late day.

The following rules apply:
\begin{itemize}[leftmargin=*]
   \item No explanation is required to use late days.
   \item At most \textbf{two late days} may be used on any single homework assignment.
   \item Late days may not be used on the final homework assignment or on any assignment for which solutions have already been released.
   \item It is the student's responsibility to keep track of how many late days have been used during the quarter.
\end{itemize}

Once all late days have been exhausted, additional late submissions will incur a penalty of \textbf{10\% per day}, up to a maximum deduction of \textbf{50\%}. Assignments submitted more than five days late, or after solutions have been released, will not be accepted.

\subsection*{Submission Issues and Technical Difficulties}

If a serious technical issue prevents a timely Gradescope submission, students may temporarily submit their assignment by email to the instructor at \href{mailto:armanjg@uw.edu}{armanjg@uw.edu}. In such cases:

\begin{itemize}[leftmargin=*]
    \item The submission must consist of a \textbf{single PDF file}. The email subject line must follow the format:
    \[
    \texttt{HW\# Submission - FULL NAME - STUDENT ID}.
    \]
    \item Once the Gradescope issue has been resolved, the student must upload the \emph{exact same file} to Gradescope as soon as possible.
    \item The student should clearly indicate in the Gradescope submission that the assignment was previously submitted by email before the deadline.
\end{itemize}

Email submissions are intended only for genuine technical emergencies and should not be used as a substitute for timely Gradescope submission.

\newpage

% ------------------------------------------------------------

\begin{problemheader}
\textbf{Problem 1.} \hfill 
\end{problemheader}
\source{23}

Three people get into an empty elevator at the first floor of a building that has 10 floors. Each presses the button for their desired floor, unless one of the others has already pressed that button. Assume that they are equally likely to want to go to floors 2 through 10, independently of each other.

What is the probability that the buttons for 3 consecutive floors are pressed?

\vfill
 
\newpage

% ------------------------------------------------------------

\begin{problemheader}
\textbf{Problem 2. } \hfill  
\end{problemheader}
\source{3}

Fred is planning to go out to dinner each night of a certain week, Monday through Friday, with each dinner being at one of his ten favorite restaurants.

\begin{enumerate}[label=(\alph*), leftmargin=*]
    \item How many possibilities are there for Fred's schedule of dinners for that Monday through Friday, if Fred is not willing to eat at the same restaurant more than once?

    \item How many possibilities are there for Fred's schedule of dinners for that Monday through Friday, if Fred is willing to eat at the same restaurant more than once, but is not willing to eat at the same place twice in a row or more?
\end{enumerate}

\vfill
\partanswerbox{(a)}
\partanswerbox{(b)}
\newpage

% ------------------------------------------------------------

\begin{problemheader}
\textbf{Problem 3.} \hfill  
\end{problemheader}
\source{15--18}

Give story proofs for the following identities.

\begin{enumerate}[label=(\alph*), leftmargin=*]
    \item Show that
    \[
    \sum_{k=0}^{n} {n \choose k}=2^n.
    \]
    

    \item Show that
    \[
    \sum_{k=0}^{n}{n \choose k}^2={2n \choose n},
    \]
    for all positive integers $n$.

    \item Show that
    \[
    \sum_{k=1}^{n} k{n \choose k}^2
    =
    n{2n-1 \choose n-1},
    \]
    for all positive integers $n$.

    \emph{Hint: Consider choosing a committee of size $n$ from two groups of size $n$ each, where only one of the two groups has people eligible to become the chair of the committee.}
\end{enumerate}

\vfill


\newpage

% ------------------------------------------------------------

\begin{problemheader}
\textbf{Problem 4. } \hfill 
\end{problemheader}
\source{25}

A city with 6 districts has 6 robberies in a particular week. Assume the robberies are located randomly, with all possibilities for which robbery occurred where equally likely.

What is the probability that some district had more than 1 robbery?

\vfill


\newpage

% ------------------------------------------------------------

\begin{problemheader}
\textbf{Problem 5. } \hfill  
\end{problemheader}
\source{5}

A knock-out tournament is being held with $2^n$ tennis players. This means that for each round, the winners move on to the next round and the losers are eliminated, until only one person remains. For example, if initially there are $2^4=16$ players, then there are 8 games in the first round, then the 8 winners move on to round 2, then the 4 winners move on to round 3, then the 2 winners move on to round 4, the winner of which is declared the winner of the tournament. There are various systems for determining who plays whom within a round, but these do not matter for this problem.

\begin{enumerate}[label=(\alph*), leftmargin=*]
    \item How many rounds are there?

    \item Count how many games in total are played, by adding up the numbers of games played in each round.

    \item Count how many games in total are played, this time by directly thinking about it without doing almost any calculation.

    \emph{Hint: How many players need to be eliminated?}
\end{enumerate}

\vfill
 
\newpage

% ------------------------------------------------------------

\begin{problemheader}
\textbf{Problem 6.} \hfill  
\end{problemheader}
\source{49}

A fair die is rolled $n$ times. What is the probability that at least 1 of the 6 values never appears?

\vfill
 
\newpage

% ------------------------------------------------------------

\begin{problemheader}
\textbf{Problem 7. } \hfill 
\end{problemheader}
\source{53}

Fred needs to choose a password for a certain website. Assume that he will choose an 8-character password, and that the legal characters are the lowercase letters $a,b,c,\dots,z$, the uppercase letters $A,B,C,\dots,Z$, and the numbers $0,1,\dots,9$.

\begin{enumerate}[label=(\alph*), leftmargin=*]
    \item How many possibilities are there if he is required to have at least one lowercase letter in his password?

    \item How many possibilities are there if he is required to have at least one lowercase letter and at least one uppercase letter in his password?

    \item How many possibilities are there if he is required to have at least one lowercase letter, at least one uppercase letter, and at least one number in his password?
\end{enumerate}

\vfill
 

\end{document}